% Фрагмент предназначен для подключения через \input.
% Требуются пакеты amsmath и amssymb.

\section{Обновление направления методом LSMR}
\label{sec:lsmr-solver}

На фиксированном внешнем шаге заданы локальные статистики
\[
 I_j\in\mathbb R^P,
 \qquad
 U_j\in\mathbb R^{P\times d},
 \qquad j\in\mathcal J.
\]
При фиксированных локальных коэффициентах \(c_j\) направление
\(\beta\in\mathbb R^d\) обновляется из ridge-задачи
\begin{equation}
 \min_{\beta\in\mathbb R^d}
 \left\{
  \sum_{j\in\mathcal J}
  \left\|I_j-c_jU_j\beta\right\|^2
  \lambda\left\|\beta-\beta_{\mathrm{prior}}\right\|^2
 \right\}.
 \label{eq:lsmr-ridge}
\end{equation}
Здесь \(\lambda=\texttt{lambda\_penalty}\), а
\(\beta_{\mathrm{prior}}\) --- направление с предыдущей внутренней итерации.

\subsection{Оператор без построения блочной матрицы}

Если формально определить
\[
 A=
 \begin{pmatrix}
  c_1U_1\\
  \vdots\\
  c_JU_J
 \end{pmatrix},
 \qquad
 b=
 \begin{pmatrix}
  I_1\\
  \vdots\\
  I_J
 \end{pmatrix},
\]
то задача \eqref{eq:lsmr-ridge} эквивалентна
\[
 \min_\beta
 \left\|
 \widetilde A\beta-\widetilde b
 \right\|^2,
 \qquad
 \widetilde A=
 \begin{pmatrix}
  A\\
  \sqrt{\lambda}\,I_d
 \end{pmatrix},
 \qquad
 \widetilde b=
 \begin{pmatrix}
  b\\
  \sqrt{\lambda}\,\beta_{\mathrm{prior}}
 \end{pmatrix}.
\]
В реализации матрицы \(A\) и \(\widetilde A\) не создаются.
Объект \texttt{LinearOperator} размера
\((JP+d)\times d\) задаёт только операции
\begin{align}
 \widetilde A v
 &=
 \begin{pmatrix}
  c_1U_1v\\
  \vdots\\
  c_JU_Jv\\
  \sqrt{\lambda}\,v
 \end{pmatrix},
 \label{eq:lsmr-matvec}\\
 \widetilde A^{\mathsf T}
 \begin{pmatrix}
  z_1\\
  \vdots\\
  z_J\\
  w
 \end{pmatrix}
 &=
 \sum_{j\in\mathcal J}c_jU_j^{\mathsf T}z_j
 \sqrt{\lambda}\,w,
 \qquad
 z_j\in\mathbb R^P,\quad w\in\mathbb R^d.
 \label{eq:lsmr-rmatvec}
\end{align}
Таким образом, LSMR решает ridge-задачу без нормальных уравнений и без
хранения полной блочной матрицы. Стоимость одной итерации составляет
\(O(JPd)\).

В коде используются параметры
\[
 \texttt{atol}=\texttt{btol}
 =\min\{\texttt{tol},10^{-8}\},
 \qquad
 \texttt{maxiter}=\max\{50,5d\}.
\]
После LSMR проверяются конечность и ненулевая норма решения, затем выполняются
нормировка и согласование знака:
\[
 \beta\leftarrow\frac{\beta}{\|\beta\|},
 \qquad
 \langle\beta,\beta_{\mathrm{prior}}\rangle<0
 \ \Longrightarrow\
 \beta\leftarrow-\beta.
\]

\subsection{Чередование с локальными коэффициентами}

Перед каждым вызовом LSMR коэффициенты пересчитываются аналитически:
\begin{equation}
 c_j
 =
 \frac{\left\langle I_j,U_j\beta\right\rangle}
 {\left\|U_j\beta\right\|^2+\eta},
 \qquad
 \eta=\texttt{local\_ridge}.
 \label{eq:lsmr-local-slopes}
\end{equation}
После этого решается \eqref{eq:lsmr-ridge}, и процесс повторяется не более
\texttt{inner\_steps} раз. Критерий остановки инвариантен к знаку направления:
\[
 \delta_\beta
 =
 \min\left\{
  \|\beta_{\mathrm{new}}-\beta_{\mathrm{prior}}\|,
  \|\beta_{\mathrm{new}}+\beta_{\mathrm{prior}}\|
 \right\}
 <\texttt{tol}.
\]
Итоговые коэффициенты \(c_j\) ещё раз вычисляются по
\eqref{eq:lsmr-local-slopes}.

\paragraph{Интерпретация остановки.}
Код \texttt{istop} метода LSMR относится только к линейной подзадаче при
фиксированных \(c_j\) и \(\beta_{\mathrm{prior}}\).
Например, \texttt{istop=2} означает, что least-squares подзадача решена,
а \texttt{istop=7} --- что LSMR достиг собственного лимита итераций.
Это не гарантирует сходимость всей чередующейся схемы: она проверяется
отдельно по \(\delta_\beta\). Кроме того, нормировка после LSMR переводит
неограниченное решение на единичную сферу, поэтому нормированный вектор уже
не является точным минимумом исходной евклидовой подзадачи
\eqref{eq:lsmr-ridge}.
